\begin{ack}

“生成”究竟是什么？在论文将尽、余绪未收之时，我依然在自问，也被追问。

上世纪七十年代末，智利科学家弗朗西斯科·瓦雷拉为认知神经科学、哲学与佛教搭起对话的桥梁。他提出的生成认知观揭示，认知具身于血肉，延展于所栖之境，并在与环境的持续耦合中涌现新的经验与觉知。若要疗救身心，必须改变经验的生活方式——以语言点燃并捕获身体的感觉，以实践校正既定的法则；刺破定见，在碎片的重组中生成新的秩序与世界。

不曾想过，这场学术的生成之旅，同样需要如此的破立重生。蜕变伴随阵痛，却也是新思考得以健康落地的必经淬炼。而这一切，承蒙多少阳光雨露的恩泽！
幸遇导师胡德文教授，他以深厚的学养与敏锐的洞察，引领我探索认知科学和人工智能交叉领域的前沿。感谢曾令李教授，使我由迷茫的求索者，逐渐成为言之有物、言必有据的思考者。感谢于扬教授，他以持续的耐心和充分的信任，支持我在探索中不断尝试与修正。亦感谢课题组其他老师在课题论证和方法设计等方面的指导与相助。及至回望，方见老师们栽下的路标；没有这些指引，我难以完成这场自我的革新。

“路过我们生命的每一个人，都参与并最终构成了我们本身”。万事由内因与外缘相缀，如树由种发芽，得水土时节，而生叶、出节、抽茎、开花、结实。至于这段经历是否值得，我想答案早已不在初始的选项里。

感谢脑机小组的程贤哲、吴安琪、李浩、童柔敏。我们携手从概念啮合到硬件落地，期待脑机头环拥有更动人的续篇。感谢李洁、张雪宁、范智鹏、刘迎欣、叶泽祺、陈子晗、李天恒、卢盖、杨俊等同门，以及潘迪博、王雨晨、张瀚文、王心雨、刘心岩、王哲等好友，让我在日常细微处感受到真诚与纯粹。或许人与人相处的美好，不完全在于炽热的表达，更是不动声色的宽容与关怀；不完全在于承诺与占有，更是彼此间自然流淌的善意与理解。

感谢黄俊文学姐。她涌流不息的乐观与生命力，洗涤着我的倦怠。我如此幸运被温柔地对待、被在意和被理解。感谢她带我接触抱石这项需要不断向上又不断摔下，但最终靠近自己的运动。它令我直面疼痛，也让我对自己诚实；它让我反复跌落，也教会我如何回到起点重新出发。

感谢家人，让我懂得，爱之为爱，在于利他。

世间万物，生住异灭。感谢一切因缘，成就当下之我；亦感激消散后的虚空，得以容纳并迎接明日之我。

\end{ack}